\section{Normalized reconstruction and feature spans}
\label{sec:reconstruction-theory}

The reconstruction formulas permit two different questions: whether a
particular map is an established kernel predictor, and whether its
coordinates give a linear decoder functions unavailable from the matched
reconstruction. We answer these questions for fixed reference geometry.
The results concern exact functions before feature standardization,
radial transformation, or penalized fitting.

Consider one local model from Section~\ref{sec:reconstruction-method} and
suppress its index. Write $M$ for its number of references,
$Q\in\R^{M\times d}$ for their row-stored coordinates, and
$C=(K+\lambda I_M)^{-1}$ for its regularized inverse, with $\lambda>0$.
The matrix $K$ is the symmetric positive-semidefinite Gaussian Gram
matrix. The vector $\mathbf1_M$ contains $M$ ones, and $k(z)\in\R^M$
contains the kernel similarities of a query $z$ to the references.
Let $A_K$, $A_R$, and $A_A$ denote the KAHM, positive-kernel RBF, and
affine reconstructions in Equation~\eqref{eq:mechanism-main-maps}.

\subsection{The meaning of normalized target transfer}

\begin{proposition}[Normalization and target transfer]
\label{prop:mechanism-krr}
For a query with $s(z)=\mathbf1_M^\top Ck(z)\ne0$, define
$\omega(z)=Ck(z)/s(z)$. Then $\mathbf1_M^\top\omega(z)=1$, and for
every vector of reference targets $y\in\R^M$,
\begin{equation}
 \omega(z)^\top y
 =\frac{k(z)^\top(K+\lambda I_M)^{-1}y}
 {k(z)^\top(K+\lambda I_M)^{-1}\mathbf1_M}.
 \label{eq:main-target-transfer}
\end{equation}
This is the ratio of two kernel-ridge predictions with regularization
matrix $K+\lambda I_M$ and no additional intercept. For an affine scalar
function $f(q)=q^\top b+c$, where $b\in\R^d$ and $c\in\R$, its
reference values $f(Q)=Qb+c\mathbf1_M$ satisfy
\[
 \omega(z)^\top f(Q)-f(z)=\{A_K(z)-z\}^\top b.
\]
\end{proposition}

The proof uses symmetry of $C$ and the fact that the normalized weights
sum to one. The numerator in Equation~\eqref{eq:main-target-transfer}
fits the supplied targets, while the denominator fits the constant
target. Because the denominator varies with the query, the ratio
generally differs from the ordinary unnormalized kernel-ridge predictor.
The affine identity also specifies the effect of reconstruction error on
an affine target: only its component in the target direction $b$ enters.
The complete proof is in Appendix~\ref{app:reconstruction-mechanism}.

\subsection{Equal spans under reference symmetry}

For a fixed common gate $g(z)$, define
$d_K(z)=g(z)\{A_K(z)-A_A(z)\}$ and
$d_R(z)=g(z)\{A_R(z)-A_A(z)\}$. Let $\mathcal S_0$ be the linear
span of the scalar gate and the coordinates of
$g(z)P^\top(z-\mu)$, where $\mu$ and $P$ are the reference center
and principal directions. The coordinates of
$g(z)A_A(z)=g(z)\mu+P\{g(z)P^\top(z-\mu)\}$ belong to
$\mathcal S_0$. Adjoining a vector-valued feature to a span means
adjoining each of its scalar coordinate functions. A common intercept
may be adjoined to both spans without changing the equality below.

\begin{proposition}[A sufficient condition for equal feature spans]
\label{prop:mechanism-redundancy}
Suppose $K\mathbf1_M=\kappa\mathbf1_M$ for a scalar $\kappa$, and
$C$ leaves
$\mathcal V=\operatorname{span}\{\mathbf1_M,Q_{:1},\ldots,Q_{:d}\}$
invariant. Here $Q_{:i}$ is column $i$ of $Q$, and invariance means
$C\mathcal V\subseteq\mathcal V$. Then
\[
 \operatorname{span}(\mathcal S_0,d_K)
 =\operatorname{span}(\mathcal S_0,d_R).
\]
More explicitly, there are $B\in\R^{d\times d}$ and $c\in\R^d$
such that, with $a=(\kappa+\lambda)^{-1}>0$,
\[
 C\mathbf1_M=a\mathbf1_M,\qquad
 CQ=QB+\mathbf1_Mc^\top,\qquad
 A_K(z)=\{B^\top A_R(z)+c\}/a.
\]
\end{proposition}

Invariance supplies the displayed representation of $CQ$. Substitution
in the normalized reconstruction expresses each KAHM departure using
RBF departures and the common affine features. The restriction of $C$
to the finite-dimensional space $\mathcal V$ is invertible onto that
space, which gives the reverse inclusion. This argument also covers
redundant reference-coordinate columns; it does not require $B$ to be
invertible. Appendix~\ref{app:reconstruction-mechanism} gives the full
proof.

A two-reference example makes the role of the common coordinates
explicit. Take scalar references $(-1,1)$, $g=1$, $P=1$, and
$A_A(z)=z$. If $\rho$ is the off-diagonal Gram entry, then
\[
 A_K=c_\rho A_R,\qquad
 d_K=c_\rho d_R+(c_\rho-1)z,\qquad
 c_\rho=\frac{1+\rho+\lambda}{1-\rho+\lambda}>0.
\]
The departures differ, but adjoining the coordinate $z$ makes their
augmented spans equal. Thus a nonlinear reconstruction formula alone
does not establish an additional raw basis.

\subsection{An irregular reference bank gives a distinct basis}

\begin{proposition}[Irregular scalar references]
\label{prop:mechanism-distinct}
Let $q_1<q_2<q_3$ satisfy $q_1+q_3\ne2q_2$, and let
$k_i(z)=\exp\{-t(z-q_i)^2/2\}$ for a precision $t>0$.
Set $K_{ij}=\exp\{-t(q_i-q_j)^2/2\}$ and
$C=(K+\lambda I_3)^{-1}$ for $\lambda>0$.
Suppose every entry of $C\mathbf1_3$ is positive. With
$Q=(q_1,q_2,q_3)^\top$, $g=1$, and $A_A(z)=z$,
\[
 A_K\notin\operatorname{span}(1,z,A_R),\qquad
 d_K\notin\operatorname{span}(1,z,d_R).
\]
These spans consist of scalar functions on the whole real line.
\end{proposition}

The six products of kernel functions have distinct exponential factors
because their exponents contain the six distinct sums $q_i+q_j$ for
$i\le j$. Their linear independence rules out an affine relation
between $A_K$ and $A_R$: such a relation would force $K$ to have
constant row sums, whereas its middle row has a strictly larger sum
than either outer row. Positivity of $C\mathbf1_3$ makes both
reconstructions bounded, excluding an additional nonzero multiple of
$z$. Appendix~\ref{app:reconstruction-mechanism} provides every step
of this argument.

For the explicit references $(-1,0,2)$, the sample-variance precision is
$t=3/7$. The reference-based regularization rule gives
$\lambda=4.498274743339644$, for which positivity of $C\mathbf1_3$
follows from the matrix bound in the appendix. This supplies a concrete
example satisfying the proposition. It does not assert distinctness
for every higher-dimensional bank fitted in the forecasting experiment.

Together, the propositions delimit the algebraic contribution. Some
reference geometries yield equal augmented spans, while an irregular
geometry can supply a distinct function. Neither conclusion determines
forecasting accuracy. A general change of feature coordinates changes
an isotropic ridge penalty, and separate nonlinear radial transformations
can remove a raw linear span relation. The empirical comparison
therefore evaluates the complete feature transformation and fitted
decoder as well as the reconstruction map.
